\scanpage{27}

\uline{Remarque}. Soient
\diagram{\begin{tikzcd}[ampersand replacement=\&]
k \arrow[r] \arrow[d,dash] \& K \arrow[d,dash] \\
k'=k\cap K^p \arrow[r,dash] \& K^p
\end{tikzcd}}
alors deux éléments $\unsure{x},y$ de $k$ peuvent être
$p$-indépendants \inserted{dans $\unsure{k|K}$ sur}
\deleted{\unsure{sur $k$}} $k\cap K^p$, et non
\inserted{dans $K$} sur $K^p$,
i.e. $d_{k/k'}x$ et $d_{k/k'}y$ peuvent être lin. indépendants,
sans que $d_{K/K^p}x$ et $d_{K/K^p}y$ le soient.
\uline{à fortiori}, $\Omega^1_{k/k\cap K^p}\otimes_k K\to\Omega^1_K$
peut n'être pas injectif, ou \deleted{\unsure{encore}} ce qui revient
au même, il peut arriver que le noyau \inserted{$N_{K/k}$}
de $\Omega_k\otimes_k K\unsure{\text{---}|\to}\Omega_K$ ne soit pas engendré sur $K$
par les $dx\otimes1$ $(x\in k\cap K^p)$.
\inserted{\uline{\unsure{qui est} \unsure{évident.} \unsure{(donc}}}

\margin{$\Omega^1_{k\cap K^p/\unsure{k^p}}\otimes_{k\cap K^p}K
\to\Omega^1_k\otimes_k\unsure{K}\to\Omega^1_K$}
\margin{N.B. \deleted{\unsure{Alors}} (\unsure{l'existence|l'assertion}
\unsure{ci-dessus}) \deleted{\unsure{à déterminer}}
\unsure{revient} \deleted{\unsure{pour}} que $N_{K/k}$
\deleted{\unsure{peut}} \unsure{n'être} ``défini sur $k$'' !
On peut \unsure{prendre} un premier \unsure{exemple} où il est
de dim \uline{1} \unsure{déjà}.}
\figure{\begin{tikzpicture}[x=1cm,y=1cm]
\draw (0,0) .. controls (.4,.5) and (.8,-.15) .. (1.2,0)
.. controls (3,.2) and (7,-.25) .. (9.3,.05)
.. controls (10,.2) and (10.8,.1) .. (11,.65);
\draw (10.85,.62)--(11,.65)--(10.94,.49);
\draw (8.2,.4) .. controls (7.65,.8) and (7.9,1.15) .. (8.5,1.18)
--(10.25,1.31);
\end{tikzpicture}}

Voici un contre. Soit $\mathbb F$ le corps premier
[ou un corps parfait quelc.], \uuline{$x,y,a_i$}
\inserted{$(1\le i\le p-1)$} des indéterminées sur $\mathbb F$,
\unsure{on pose} \deleted{$\unsure{\beta=y-\alpha x}$}
\[
\deleted{K^p=\mathbb F(x^p,\unsure{y^p},\unsure{\alpha})
\qquad\text{donc}\quad K=\mathbb F(x,\unsure{y},\unsure{\alpha^{1/p}})}
\]
\deleted{$\unsure{k'}=$}
\[
\left\{\begin{aligned}
a_0&=y-\sum_1^{p-1}a_i x^i\\
a&=x^p
\end{aligned}\right.
\qquad \text{soit}\quad y=\sum_0^{p-1}a_i x^i
\]
\[
K^p=\mathbb F(a,a_0,\unsure{a_1,\ldots,a_{p-1}})
=\deleted{\mathbb F(\unsure{\alpha},x^p,y-\alpha x)}
\inserted{\mathbb F\Bigl(x^p,a_1,\ldots,a_{p-1},y-\sum_1^{p-1}a_i x^i\Bigr)}
\]
\[
\deleted{K=\mathbb F(\alpha^{1/p},x,(y-\alpha x)^{1/p})}
\qquad\text{i.e.}\quad K=\mathbb F(x,a_1^{1/p},\ldots,a_{p-1}^{1/p},(\text{''})^{1/p})
\]
\[
k=\mathbb F(x,y)
\]
\margin{NB $a$ et $(a_i)_{0\le i\le p-1}$ sont alg. ind. sur $\mathbb F$.}
\figure{\begin{tikzpicture}[x=1cm,y=1cm]
\draw (2.8,1.2) .. controls (3.1,1.2) and (2.9,.85) .. (3.12,.8)
.. controls (3.3,.75) and (3.1,.5) .. (3.2,.3)--(2.9,.25);
\draw (.9,.08)--(2.65,.35);
\end{tikzpicture}}
\margin{[N.B. $x,\unsure{z}\in K$, donc $y=\unsure{\alpha x+\beta}\in K$,
donc $\unsure{k}\subset K$}

\mspage{26}
